% frontmatter/polish/actividades.tex -- las actividades, una a una, para
% el adulto (frontmatter/actividades.tex, en polaco): esta página también
% tiene que caber en una (tools/check_pages.py comprueba que el día 1 está
% en la página 5).
\section*{Zadania}

Każde zadanie pojawia się w dniu, w którym jest potrzebne po raz
pierwszy, a potem co jakiś czas wraca. Polecenia na każdej stronie, małą
szarą czcionką, może przeczytać dorosły, jeśli trzeba.

{\small
\begin{multicols}{2}
\raggedright
\subsection*{Od jesieni}
\begin{itemize}[leftmargin=5mm, itemsep=1pt, topsep=0pt]
\item \textbf{\lblClasifica}: przeczytać każdą karteczkę i połączyć ją z
  jej pudełkiem -- na przykład \emph{potrzebuję} albo \emph{chcę}.
\item \textbf{\lblUne}: połączyć każdą rzecz po lewej z tą, która do
  niej pasuje, po prawej.
\item \textbf{\lblVF}: przeczytać każde zdanie i zakreślić P, jeśli jest
  prawdziwe, albo F, jeśli jest fałszywe.
\item \textbf{\lblTrueque}: jeśli tyle rzeczy można zamienić na tyle
  innych, ile dostanie się za te? Można je narysować.
\item \textbf{\lblDibuja}: narysować to, o czym mowa.
\item \textbf{\lblRepasa} (w piątki): słowo tygodnia, to, co dziecko już
  umie, do zaznaczenia, i rysunek. Co dziesięć stron jest mały napis, żeby
  to uczcić.
\item \textbf{\lblDinero}: policzyć najpierw banknoty, a potem monety, i
  wpisać, ile jest euro.
\item \textbf{\lblReparte}: dzielić po jednym, rysując na każdym talerzu
  (albo w każdej ramce) to, co dostaje każdy, i zobaczyć, co zostaje.
\item \textbf{\lblCompra}: znaleźć cenę każdej kupowanej rzeczy i dodać
  je.
\item \textbf{\lblProblema}: krótkie zadanie z działaniem; można je
  narysować w ramce.
\item \textbf{\lblBotes}: pieniądze w trzech słoikach -- na
  oszczędności, na wydatki i do dzielenia się --; w jednym brakuje tego,
  co w nim jest.
\item \textbf{\lblHucha}: co tydzień odkładać tyle samo, zapisując, ile
  już jest, aż do celu.
\item \textbf{\lblLlega}: zakreślić to, co można kupić za pieniądze,
  które się ma.
\item \textbf{\lblVuelta}: reszta to to, co się daje, minus to, ile
  rzecz kosztuje.
\item \textbf{\lblOrdena}: ułożyć kroki po kolei, wpisując numer w każdą
  kratkę.
\end{itemize}
\subsection*{Od zimy}
\begin{itemize}[leftmargin=5mm, itemsep=1pt, topsep=0pt]
\item \textbf{\lblElige}: pieniędzy wystarcza tylko na jedną rzecz:
  zakreślić tę, którą się wybiera, skreślić tę, z której się rezygnuje, i
  zobaczyć, ile zostaje. Każdy wybór jest dobry.
\item \textbf{\lblCompara}: ta sama rzecz w dwóch sklepach: zakreślić
  sklep, w którym jest taniej, i wpisać, ile się oszczędza.
\end{itemize}
\subsection*{Od wiosny}
\begin{itemize}[leftmargin=5mm, itemsep=1pt, topsep=0pt]
\item \textbf{\lblCuentas}: zeszyt rachunków: to, co przychodzi, się
  dodaje, to, co wychodzi, się odejmuje, a w każdą kratkę wpisuje się,
  ile zostaje.
\end{itemize}
\end{multicols}}
\newpage
